\documentclass[14pt]{extarticle}
\usepackage{fontspec}
\usepackage{polyglossia}
\usepackage{amsmath}
\usepackage{amssymb}
\usepackage{xcolor}
\usepackage{fancyhdr}
\usepackage{graphicx}
\usepackage{listings}
\usepackage[most]{tcolorbox}
% \usepackage{setspace}
\usepackage{pifont}
\usepackage{enumitem}
\usepackage{titlesec}
\usepackage[bottom]{footmisc}
\usepackage{titling}
\usepackage{minted}
\usepackage{etoolbox}
\usepackage{array}
\usepackage{longtable} % page-breakable tables (must load before polyglossia/bidi RTL setup)
\usepackage{extsizes}
\usepackage{float}
\usepackage{booktabs}
\usepackage{hyperref}
\usepackage[onehalfspacing]{setspace}
\newif\ifwhiteblack
\ifdefined\setwhiteblack
  \setwhiteblack
\fi

\ifwhiteblack
  \usepackage[a4paper, landscape, top=1.25cm, bottom=1.25cm, left=1.25cm, right=1.25cm, headheight=1.5cm, headsep=0.5cm]{geometry}
\else
  \usepackage[a4paper, top=2.7cm, bottom=2.7cm, left=2.5cm, right=2.5cm, headheight=1.5cm, headsep=0.5cm]{geometry}
\fi
\usepackage{tikz}
\usepackage{datetime2}

\usetikzlibrary{automata, positioning, arrows.meta, arrows}
\providecommand{\LTcaptype}{table}

\newfontfamily\emoji{Segoe UI Emoji}

\pagestyle{fancy}

\setmainlanguage[numerals=western]{arabic}
\setotherlanguage{english}
\newfontfamily\arabicfont[Script=Arabic]{Calibri}
\newfontfamily\arabicfonttt[Script=Arabic]{Courier New}



\ifwhiteblack
\lstset{
  language=[Sharp]C,
  numbers=left,
  stepnumber=1,
  numbersep=8pt,
  frame=single,
  basicstyle=\ttfamily\small,
  keywordstyle=\bfseries,
  stringstyle=\ttfamily,
  commentstyle=\itshape
}
\else
\lstset{
  language=[Sharp]C,
  numbers=left,
  stepnumber=1,
  numbersep=8pt,
  frame=single,
  basicstyle=\ttfamily\small,
  keywordstyle=\color{blue},
  stringstyle=\color{red},
  commentstyle=\color{green!50!black}
}
\fi

\newif\ifdetailed
\ifdefined\setdetailed
  \setdetailed
\fi

\newif\ifwithsols
\ifdefined\setwithsols
  \setwithsols
\fi

\ifwhiteblack
  \AtBeginDocument{
    \LARGE
    \usemintedstyle{bw}
  }
\fi

% Default minted look: compact code that wraps long lines.
% Not applied in black-and-white landscape mode, which already uses \LARGE.
\ifwhiteblack\else
  \setminted{fontsize=\small, breaklines=true}
\fi

% unified tcolorboxes for articles
\ifwhiteblack
\tcbset{colback=white, colframe=black, fonttitle=\bfseries, boxrule=1pt}
\newtcolorbox{boxDef}[1][]{title={{\emoji📘} تعريف\ifx\\#1\\\else ~#1\fi :}}
\newtcolorbox{boxExercise}[1][]{title={{\emoji🧩} تمرين\ifx\\#1\\\else ~#1\fi :}}
\newtcolorbox{boxExample}[1][]{title={{\emoji📝} مثال\ifx\\#1\\\else ~#1\fi :}}
\newtcolorbox{boxNote}[1][]{title={{\emoji✨} ملاحظة\ifx\\#1\\\else ~#1\fi :}}
\newtcolorbox{boxAttention}[1][]{title={{\emoji🔔} تنبيه\ifx\\#1\\\else ~#1\fi :}}
\newtcolorbox{boxWarning}[1][]{title={{\emoji⚡} ملاحظة هامة\ifx\\#1\\\else ~#1\fi :}}
\newtcolorbox{boxSolution}[1][]{title={{\emoji✅} حل\ifx\\#1\\\else ~#1\fi :}}
\newtcolorbox{boxSymbol}[1][]{title={{\emoji🔣} رمز\ifx\\#1\\\else ~#1\fi :}}
\newtcolorbox{boxHint}[1][]{title={{\emoji💡} تلميح\ifx\\#1\\\else ~#1\fi :}}
\else
\tcbset{colback=white, colframe=black, fonttitle=\bfseries, boxrule=0.8pt}
\newtcolorbox{boxDef}[1][]{colback=blue!5!white,colframe=blue!75!black,
  title={{\emoji📘} تعريف\ifx\\#1\\\else ~#1\fi :}}
\newtcolorbox{boxExercise}[1][]{colback=cyan!5!white,colframe=cyan!70!black,
  title={{\emoji🧩} تمرين\ifx\\#1\\\else ~#1\fi :}}
\newtcolorbox{boxExample}[1][]{colback=yellow!5!white,colframe=orange!90!black,
  title={{\emoji📝} مثال\ifx\\#1\\\else ~#1\fi :}}
\newtcolorbox{boxNote}[1][]{colback=gray!10!white,colframe=black,
  title={{\emoji✨} ملاحظة\ifx\\#1\\\else ~#1\fi :}}
\newtcolorbox{boxAttention}[1][]{colback=magenta!10!white,colframe=magenta!80!black,
  title={{\emoji🔔} تنبيه\ifx\\#1\\\else ~#1\fi :}}
\newtcolorbox{boxWarning}[1][]{colback=red!5!white,colframe=red!75!black,
  title={{\emoji⚡} ملاحظة هامة\ifx\\#1\\\else ~#1\fi :}}
\newtcolorbox{boxSolution}[1][]{colback=green!5!white,colframe=green!60!black,
  title={{\emoji✅} حل\ifx\\#1\\\else ~#1\fi :}}
\newtcolorbox{boxSymbol}[1][]{colback=purple!5!white,colframe=purple!70!black,
  title={{\emoji🔣} رمز\ifx\\#1\\\else ~#1\fi :}}
\newtcolorbox{boxHint}[1][]{colback=teal!5!white,colframe=teal!60!black,
  title={{\emoji💡} تلميح\ifx\\#1\\\else ~#1\fi :}}
\fi


\ifwhiteblack
\tcbset{simplecode/.style={ colback=white, colframe=black, boxrule=1pt, arc=2pt, left=4pt,right=4pt,top=4pt,bottom=4pt}}
\else
\tcbset{simplecode/.style={ colback=gray!5, colframe=black!50, boxrule=0.4pt, arc=2pt, left=4pt,right=4pt,top=4pt,bottom=4pt}}
\fi
% Compact code boxes: tighter line spacing (the preamble uses onehalfspacing) and small vertical gaps.
% Must come AFTER the simplecode definition above.
\tcbset{simplecode/.append style={fontupper=\setstretch{0.95}, top=4pt, bottom=4pt, before skip=6pt, after skip=6pt}}
\newenvironment{boxCode}{\begin{tcolorbox}[simplecode]}{\end{tcolorbox}}

% Code cell used inside tables (keeps code left-to-right inside RTL tables).
% Use inside a fixed-width p{...} column: \begin{pycell}\begin{minted}...\end{minted}\end{pycell}
\newtcolorbox{pycell}{frame empty,colback=white,boxsep=0pt,left=2pt,right=2pt,top=2pt,bottom=2pt,width=\linewidth}

\newcolumntype{C}[1]{>{\centering\arraybackslash}p{#1}}

\makeatletter
\DTMnewdatestyle{dotted}{%
  \renewcommand{\DTMdisplaydate}[4]{\two@digits{##3}.\two@digits{##2}.##1}%
  \renewcommand{\DTMDisplaydate}{\DTMdisplaydate}%
}
\makeatother
\DTMsetdatestyle{dotted}

% redefine spaces after titles
\makeatletter
\renewcommand{\@maketitle}{%
  \begin{center}
    {\huge \bfseries \@title \par}%
    \vskip 0.2em % space between title and author
    {\large \@author \par}%
    % \vskip 0.2em % space between author and date
    % {\normalsize \@date \par}%
    {\small \hfill \begin{english}\DTMtoday\end{english} \par}%
  \end{center}
}
\makeatother

\fancyhf{} % clear default
\fancypagestyle{plain}{
  \fancyhf{}
  \fancyhead[L]{مدرسة التسامح الشاملة}
  % \fancyhead[L]{\includegraphics[height=1cm]{../../../images/logoTasamoh.png}}
  \fancyhead[R]{الأستاذ محمود اغبارية}
  \fancyfoot[C]{\thepage}
}

\fancyhead[L]{مدرسة التسامح الشاملة}
\fancyhead[R]{الأستاذ محمود اغبارية}
\fancyfoot[C]{\thepage}

\setcounter{tocdepth}{3} % only section subsection and subsubsection in TOC

\newcommand{\importpdfpage}[2]{%
    \thispagestyle{empty}
    \newgeometry{left=0cm,right=0cm,top=0cm,bottom=0cm}
    \includegraphics[page=#2,width=0.95\paperwidth,height=0.95\paperheight,keepaspectratio]{#1}
    \restoregeometry
}

\newcommand{\insertFullImg}[1]{%
    \noindent
    \makebox[\textwidth][c]{
        \includegraphics[width=0.9\paperwidth,keepaspectratio]{#1}%
    }%
}%


\title{ملخص: العمليات الخارجية في لغة \textenglish{C\#}}

\begin{document}
\maketitle
\thispagestyle{fancy}

\section*{مقدمة}
العمليات الخارجية هي وسيلة لتقسيم البرنامج إلى مقاطع صغيرة ومستقلة تؤدي مهام محددة.
تساعد العمليات على:
\begin{itemize}
    \item تنظيم الكود وجعله أسهل للقراءة والفهم والصيانة.
    \item إعادة استخدام نفس الكود أكثر من مرة دون تكرار كتابته.
\end{itemize}

في هذا الدرس سنتعرّف على مبنى العملية الخارجية، والفرق بين تلقي المعطيات واستقبال المدخلات، وأقسام العمليات المختلفة وكيفية تعريفها واستدعائها داخل الفئة.

\begin{boxAttention}[- مصطلحات أساسية مهمة جدًّا]
\begin{itemize}
    \item \textbf{تستقبل:} تقرأ قيمة يُدخلها المستخدم عبر لوحة المفاتيح باستخدام \textenglish{Console.ReadLine()}.
    \item \textbf{تتلقى:} تأخذ معطيات (بارامترات) تُمرَّر إليها بين القوسين في تعريف العملية وعند استدعائها.
    \item \textbf{تطبع:} تعرض مخرجات على الشاشة باستخدام \textenglish{Console.WriteLine(...)}.
    \item \textbf{تعيد:} تُرجع قيمة إلى المكان الذي استُدعيت منه باستخدام الأمر \textenglish{return}.
\end{itemize}
\end{boxAttention}

\clearpage
% ==============================================================
% 1. مبنى تعريف العملية واستدعاؤها
% ==============================================================
\section{مبنى تعريف العملية واستدعاؤها}

كل عملية خارجية في لغة \textenglish{C\#} تتكوّن من قسمين رئيسيين:

\begin{enumerate}
    \item \textbf{تعريف العملية:}
    وهو السطر الأوّل الذي يحتوي على:
    \begin{itemize}
        \item الكلمتين الثابتتين في هذه المرحلة: \textenglish{public static}
        \item نمط القيمة التي تعيدها العملية (\textenglish{int, double, bool, string, char})، أو \textenglish{void} إذا كانت لا تعيد أي قيمة.
        \item اسم العملية، ومن المتّبع أن يبدأ دائمًا بحرف كبير (\textenglish{PascalCase}).
        \item المعطيات (البارامترات) التي تتلقاها العملية بين القوسين، مثل \textenglish{(int x, int y)}، أو قوسين فارغين \textenglish{()} إذا كانت لا تتلقى معطيات.
    \end{itemize}

    \item \textbf{جسم العملية:}
    وهو الكود الموجود بين القوسين المتعرجين \textenglish{\{ \}}، ويحتوي على التعليمات التي تُنفّذ عند استدعاء العملية.
\end{enumerate}

\filbreak
\begin{boxExample}[1: عملية لا تتلقى ولا تعيد]
\begin{english}
\begin{minted}{csharp}
public static void PrintHello()
{
    Console.WriteLine("Hello");
}
\end{minted}
\end{english}
\hrule
\vspace{0.4em}
\begin{itemize}[nosep]
    \item \textenglish{public static} $\leftarrow$ جزء ثابت في تعريف العملية الخارجية.
    \item \textenglish{void} $\leftarrow$ العملية لا تعيد أي قيمة.
    \item \textenglish{PrintHello} $\leftarrow$ اسم العملية.
    \item \textenglish{()} $\leftarrow$ الأقواس الفارغة بعد اسم العملية تعني أنّ العملية لا تتلقى معطيات (بارامترات).
    \item ما بين القوسين \textenglish{\{ \}} هو جسم العملية، وفي هذه الحالة يطبع الرسالة \textenglish{Hello}.
\end{itemize}
\end{boxExample}

\filbreak
\begin{boxExample}[2: عملية تتلقى وتعيد]
\begin{english}
\begin{minted}{csharp}
public static int AddNumbers(int a, int b)
{
    int sum = a + b;
    return sum;
}
\end{minted}
\end{english}
\hrule
\vspace{0.4em}
\begin{itemize}[nosep]
    \item \textenglish{public static} $\leftarrow$ جزء ثابت.
    \item \textenglish{int} $\leftarrow$ نمط القيمة المعادة.
    \item \textenglish{AddNumbers} $\leftarrow$ اسم العملية.
    \item \textenglish{(int a, int b)} $\leftarrow$ المعطيات (البارامترات) التي تتلقاها العملية، وفي هذه الحالة تتلقى متغيرين كلاهما من نمط \textenglish{int}.
    \item نستخدم الأمر \textenglish{return} لإنهاء العملية وإعادة القيمة المطلوبة. \\
    \textbf{ملاحظة:} نمط القيمة التي تعيدها العملية في تعريفها يجب أن يطابق نمط القيمة المكتوبة بعد كلمة \textenglish{return}.
\end{itemize}
\end{boxExample}

\filbreak
\subsection*{كيفية استدعاء العملية في \textenglish{Main}}
عندما تكون العملية معرّفة داخل نفس الفئة (\textenglish{Program})، نستدعيها بكتابة اسمها مباشرة:

\begin{boxExample}[استدعاء العمليات في \textenglish{Main}]
\begin{english}
\begin{minted}{csharp}
public static void Main(string[] args)
{
    PrintHello();
    int sum1 = AddNumbers(5, 10);
    int a = 5;
    int b = 10;
    int sum2 = AddNumbers(a, b);
}
\end{minted}
\end{english}
\end{boxExample}

\begin{itemize}
    \item العملية \textenglish{PrintHello()} لا تعيد أي قيمة (\textenglish{void})، لذلك نكتب اسمها فقط ولا نحفظها في متغير.
    \item العملية \textenglish{AddNumbers(int a, int b)} تعيد قيمة من نمط \textenglish{int}، لذلك نحفظ القيمة التي تعيدها في متغير من نفس النمط (\textenglish{sum1} أو \textenglish{sum2})، أو نطبعها مباشرة.
    \item المعطيات التي نمرّرها للعملية عند استدعائها يجب أن تطابق المعطيات في تعريف العملية من حيث العدد والنمط والترتيب.
\end{itemize}

\filbreak
\subsection*{استدعاء عملية خارجية من فئة أخرى}
لكي نستدعي عملية معرّفة في فئة أخرى، نكتب اسم الفئة ثم نقطة \textenglish{.} ثم اسم العملية. نحن نعرف ونستخدم الكثير من هذه العمليات الجاهزة:
\begin{itemize}[nosep]
    \item \textenglish{Console.WriteLine("Hello");} $\leftarrow$ استدعاء للعملية \textenglish{WriteLine} من الفئة \textenglish{Console} (تتلقى نصًّا من نمط \textenglish{string} ولا تعيد قيمة).
    \item \textenglish{Console.ReadLine();} $\leftarrow$ استدعاء للعملية \textenglish{ReadLine} من الفئة \textenglish{Console} (لا تتلقى معطيات وتعيد قيمة من نمط \textenglish{string}).
    \item \textenglish{int.Parse(str);} و \textenglish{double.Parse(str);} $\leftarrow$ تتلقى كل منهما نصًّا من نمط \textenglish{string} وتعيد العدد المقابل (\textenglish{int} أو \textenglish{double}).
    \item عمليات الفئة \textenglish{Math}: مثل \textenglish{Math.Max(5, 10)}، \textenglish{Math.Min(5, 10)}، \textenglish{Math.Abs(-5)}، \textenglish{Math.Pow(2, 3)}، \textenglish{Math.Sqrt(25)}، و \textenglish{Math.Round(2.5)}.
\end{itemize}

\clearpage
% ==============================================================
% 2. أقسام العمليات الخارجية (حسب تقسيم الأسئلة)
% ==============================================================
\section{أقسام العمليات الخارجية}

تُقسّم العمليات الخارجية بحسب تلقي المعطيات ونمط القيمة المعادة إلى الأقسام التالية:

% --------------------------------------------------------------
% القسم الأول
% --------------------------------------------------------------
\filbreak
\subsection{القسم الأوّل: عمليات لا تتلقى ولا تعيد (تستقبل وتطبع فقط)}
في هذا القسم يكون نمط الإرجاع \textenglish{void} والأقواس بعد اسم العملية فارغة \textenglish{()}. تقوم العملية باستقبال المدخلات من المستخدم عبر \textenglish{Console.ReadLine()} وطباعة المطلوب مباشرة.

\begin{boxExample}[عملية لا تتلقى ولا تعيد - تستقبل عددين وتطبع حاصل ضربهما]
\begin{english}
\begin{minted}{csharp}
public static void ReadAndPrintProduct()
{
    Console.WriteLine("Enter two integers:");
    int a = int.Parse(Console.ReadLine());
    int b = int.Parse(Console.ReadLine());
    Console.WriteLine("Product = " + (a * b));
}

public static void Main()
{
    ReadAndPrintProduct();
}
\end{minted}
\end{english}
\end{boxExample}
\begin{itemize}[nosep]
    \item لا تتلقى $\leftarrow$ نستدعيها بدون معطيات بين القوسين: \textenglish{ReadAndPrintProduct();}.
    \item لا تعيد (\textenglish{void}) $\leftarrow$ لا يوجد فيها \textenglish{return}، ولا نحفظ استدعاءها في متغير.
\end{itemize}

% --------------------------------------------------------------
% القسم الثاني
% --------------------------------------------------------------
\filbreak
\subsection{القسم الثاني: عمليات تتلقى ولا تعيد (قسم منها يستقبل أيضًا)}
تتلقى العملية معطى واحدًا أو أكثر بين القوسين، وتنفّذ مهمتها وتطبع الناتج دون إرجاع قيمة (\textenglish{void}). ويمكن للعملية أن تستقبل مدخلات إضافية من المستخدم إلى جانب المعطيات التي تلقتها.

\begin{boxExample}[1: عملية تتلقى عددين وتطبع مجموعهما]
\begin{english}
\begin{minted}{csharp}
public static void PrintSum(int a, int b)
{
    int sum = a + b;
    Console.WriteLine("Sum = " + sum);
}

public static void Main()
{
    PrintSum(5, 3);
    int x = 10;
    int y = 20;
    PrintSum(x, y);
}
\end{minted}
\end{english}
\end{boxExample}

\filbreak
\begin{boxExample}[2: عملية تتلقى أجر الساعة وتستقبل عدد ساعات العمل ثم تطبع الراتب]
\begin{english}
\begin{minted}{csharp}
public static void PrintSalary(double hourlyRate)
{
    Console.WriteLine("Enter hours worked:");
    int hours = int.Parse(Console.ReadLine());
    double salary = hourlyRate * hours;
    Console.WriteLine("Salary = " + salary);
}

public static void Main()
{
    PrintSalary(45.0);
}
\end{minted}
\end{english}
\end{boxExample}
\begin{itemize}[nosep]
    \item تتلقى $\leftarrow$ نمرّر لها عند الاستدعاء معطيات من نفس النمط المحدد في تعريفها.
    \item لا تعيد (\textenglish{void}) $\leftarrow$ تطبع النتيجة داخلها ولا تعيد قيمة إلى \textenglish{Main}.
\end{itemize}

% --------------------------------------------------------------
% القسم الثالث
% --------------------------------------------------------------
\filbreak
\subsection{القسم الثالث: عمليات لا تتلقى وتعيد أعدادًا صحيحة أو عشرية}
أقواس العملية فارغة \textenglish{()} لأنّها لا تتلقى معطيات، لكنّها تعيد عددًا (\textenglish{int} أو \textenglish{double}) باستخدام \textenglish{return}. غالبًا تستقبل هذه العملية قيمًا من المستخدم، وقد تطبع معلومات إضافية قبل إرجاع النتيجة.

\begin{boxExample}[1: عملية تستقبل عددًا من المستخدم وتعيده]
\begin{english}
\begin{minted}{csharp}
public static int ReadNumber()
{
    Console.WriteLine("Enter an integer:");
    int num = int.Parse(Console.ReadLine());
    return num;
}

public static void Main()
{
    int x = ReadNumber();
    Console.WriteLine("The number is: " + x);
}
\end{minted}
\end{english}
\end{boxExample}

\filbreak
\begin{boxExample}[2: عملية تستقبل سعرين، تطبع مجموعهما، وتعيد قيمة الضريبة ($17\%$)]
\begin{english}
\begin{minted}{csharp}
public static double ReadPricesPrintSumGetTax()
{
    double p1 = double.Parse(Console.ReadLine());
    double p2 = double.Parse(Console.ReadLine());
    double sum = p1 + p2;
    Console.WriteLine("Sum before tax = " + sum);
    return sum * 0.17;
}

public static void Main()
{
    double tax = ReadPricesPrintSumGetTax();
    Console.WriteLine("Tax = " + tax);
}
\end{minted}
\end{english}
\end{boxExample}

% --------------------------------------------------------------
% القسم الرابع
% --------------------------------------------------------------
\filbreak
\subsection{القسم الرابع: عمليات تتلقى وتعيد أعدادًا صحيحة أو عشرية}
تتلقى العملية معطيات بين القوسين، وتجري حسابات عليها (وقد تستقبل مدخلات إضافية أو تطبع رسائل)، ثم تعيد الناتج العددي (\textenglish{int} أو \textenglish{double}) باستخدام \textenglish{return}.

\begin{boxExample}[1: عملية تتلقى طول مستطيل وعرضه وتعيد مساحته]
\begin{english}
\begin{minted}{csharp}
public static double GetRectangleArea(double length, double width)
{
    double area = length * width;
    return area;
}

public static void Main()
{
    double result = GetRectangleArea(6.5, 4.0);
    Console.WriteLine("Area = " + result);
}
\end{minted}
\end{english}
\end{boxExample}

\filbreak
\begin{boxExample}[2: عملية تتلقى الرصيد الحالي، وتستقبل مبلغ الإيداع، وتطبع رسالة وتعيد الرصيد الجديد]
\begin{english}
\begin{minted}{csharp}
public static double DepositMoney(double balance)
{
    Console.WriteLine("Enter deposit amount:");
    double amount = double.Parse(Console.ReadLine());
    Console.WriteLine("Deposit completed");
    return balance + amount;
}

public static void Main()
{
    double newBalance = DepositMoney(500.0);
    Console.WriteLine("New balance = " + newBalance);
}
\end{minted}
\end{english}
\end{boxExample}

% --------------------------------------------------------------
% القسم الخامس
% --------------------------------------------------------------
\filbreak
\subsection{القسم الخامس: عمليات تعيد \textenglish{string} أو \textenglish{char}}
يمكن للعملية أن تعيد نصًّا (\textenglish{string}) أو حرفًا واحدًا (\textenglish{char})، سواء كانت تتلقى معطيات، أو تستقبل مدخلات من المستخدم، أو لا تتلقى ولا تستقبل.

\begin{boxExample}[1: عملية تتلقى اسمًا وتعيد نصًّا ترحيبيًّا (\textenglish{string})]
\begin{english}
\begin{minted}{csharp}
public static string Greet(string name)
{
    return "Hello " + name + "!";
}

public static void Main()
{
    string msg = Greet("Ali");
    Console.WriteLine(msg);
}
\end{minted}
\end{english}
\end{boxExample}

\filbreak
\begin{boxExample}[2: عملية تتلقى عددًا صحيحًا وتعيد رمز إشارته (\textenglish{char})]
\begin{english}
\begin{minted}{csharp}
public static char GetSignSymbol(int num)
{
    if (num >= 0)
    {
        return '+';
    }
    else
    {
        return '-';
    }
}

public static void Main()
{
    char symbol = GetSignSymbol(-12);
    Console.WriteLine("Sign: " + symbol);
}
\end{minted}
\end{english}
\end{boxExample}

% --------------------------------------------------------------
% القسم السادس
% --------------------------------------------------------------
\clearpage
\subsection{القسم السادس: عمليات تعيد \textenglish{bool} وكيفية استخدامها داخل جملة \textenglish{if}}
العمليات التي تعيد قيمة منطقية (\textenglish{bool}) تفحص شرطًا معيّنًا وتعيد \textenglish{true} إذا تحقّق الشرط، أو \textenglish{false} إذا لم يتحقّق. \\
بما أنّ جملة \textenglish{if} تتوقّع بين قوسيها تعبيرًا منطقيًّا قيمته \textenglish{true} أو \textenglish{false}، يمكننا \textbf{استدعاء العملية التي تعيد \textenglish{bool} مباشرة داخل قوسي \textenglish{if}}.

\begin{boxExample}[1: تعريف عملية تعيد \textenglish{bool} واستدعاؤها مباشرة داخل \textenglish{if}]
\begin{english}
\begin{minted}{csharp}
public static bool IsPositive(int num)
{
    if (num > 0)
    {
        return true;
    }
    else
    {
        return false;
    }
}

public static void Main()
{
    int n = int.Parse(Console.ReadLine());
    if (IsPositive(n))
    {
        Console.WriteLine("Positive number");
    }
    else
    {
        Console.WriteLine("Zero or negative");
    }
}
\end{minted}
\end{english}
\end{boxExample}

\filbreak
\begin{boxExample}[2: كتابة عملية تعيد \textenglish{bool} باختصار واستخدامها في \textenglish{if}]
\begin{english}
\begin{minted}{csharp}
public static bool IsAdult(int age)
{
    return (age >= 18);
}

public static void Main()
{
    int userAge = int.Parse(Console.ReadLine());
    if (IsAdult(userAge))
    {
        Console.WriteLine("Allowed");
    }
    else
    {
        Console.WriteLine("Under 18");
    }
}
\end{minted}
\end{english}
\end{boxExample}

\begin{itemize}[nosep]
    \item عند تنفيذ \textenglish{if (IsPositive(n))}، يتم أولاً استدعاء العملية \textenglish{IsPositive(n)}؛ فإذا أعادت \textenglish{true} يُنفَّذ كود \textenglish{if}، وإذا أعادت \textenglish{false} يُنفَّذ كود \textenglish{else}.
    \item يمكن أيضًا حفظ النتيجة في متغير منطقي أولاً ثم فحصه: \textenglish{bool res = IsPositive(n); if (res) \{ ... \}}.
\end{itemize}

% --------------------------------------------------------------
% القسم السابع
% --------------------------------------------------------------
\filbreak
\subsection{القسم السابع: عمليات تستدعي عمليات أخرى}
تمامًا كما نستدعي العمليات الخارجية من داخل \textenglish{Main}، يمكن لأي عملية خارجية أن تستدعي عملية خارجية أخرى لمساعدتها في أداء مهمتها.

\begin{boxExample}[عملية تحسب تكعيب عدد باستدعاء عملية تضرب عددين]
\begin{english}
\begin{minted}{csharp}
public static int Multiply(int a, int b)
{
    return a * b;
}

public static int Cube(int x)
{
    int temp = Multiply(x, x);
    return Multiply(temp, x);
}

public static void Main()
{
    int res = Cube(4);
    Console.WriteLine("Result = " + res);
}
\end{minted}
\end{english}
\end{boxExample}

\filbreak
\section{ملاحظات هامة}
\begin{itemize}
    \item يُفضّل أن يبدأ اسم العملية بحرف كبير وأن يصف وظيفتها بوضوح، وإذا كان مكوّنًا من أكثر من كلمة فكل كلمة تبدأ بحرف كبير (\textenglish{CalculateSum}، \textenglish{GetMaxValue}، \textenglish{PrintResult}).
    \item اجعل كل عملية تقوم بمهمة واحدة واضحة، واستخدم العمليات بدلاً من تكرار الكود.
    \item داخل كل عملية يمكنك كتابة جميع الأوامر التي تعلمتها (تعريف متغيرات، استقبال، طباعة، جمل شرطية \textenglish{if - else}، واستدعاء عمليات أخرى).
    \item يمكن تمرير القيمة التي تعيدها عملية كمعطى لعملية أخرى مباشرة، مثل: \textenglish{Math.Max(a, Math.Max(b, c))}.
\end{itemize}

\vspace{1cm}
\begin{flushleft}
أرجو لكم وقتًا ممتعًا.

الأستاذ محمود اغبارية.
\end{flushleft}

\end{document}
